\section{Система исполнения инструкций}

\subsection{Таблица диспетчеризации}

\subsubsection{Структура: фабрика и виртуальный вызов}

Каждая PTX-инструкция представлена отдельным C++-классом,
наследующим абстрактный базовый класс \texttt{Instruction}.
Создание объекта нужного подкласса выполняется фабричной функцией \\
\texttt{makeInstruction(name, line)}, которую генерирует
\texttt{autogen.py} из \\\texttt{instructions.json}:

\begin{lstlisting}[language=C++, caption={Фрагмент автогенерируемой фабрики}]
std::shared_ptr<Instruction> makeInstruction(
        const std::string& name, const std::string& line) {
    if (name == "add")  return addInstruction::Make(line);
    if (name == "fma")  return fmaInstruction::Make(line);
    if (name == "ld")   return ldInstruction::Make(line);
    // ... all supported instructions ...
}
\end{lstlisting}

Фабрика вызывается один раз при разборе PTX-текста. В горячем цикле
исполнения варпа производится единственный виртуальный вызов:
\texttt{instr->Execute(warp\_ctx)}.
Накладные расходы~--- одна косвенная адресация через таблицу виртуальных методов (vtable).

\subsubsection{Генерация кода классов инструкций}

Скрипт \texttt{autogen.py} читает \texttt{instructions.json}
и генерирует файлы \texttt{instructions.h}/\texttt{.cpp}
с классом для каждой мнемоники. Единственный источник истины
в JSON предотвращает рассинхронизацию между парсером, фабрикой
и тестами.

Принцип открытости/закрытости: добавление поддержки новой инструкции
сводится к добавлению записи
в \texttt{instructions.json} и реализации метода \texttt{Execute()}~---
\texttt{autogen.py} вставит нужную ветку в фабрику автоматически.

\subsection{Регистровый файл и разрешение операндов}

\subsubsection{Структура регистрового файла}

Регистровый файл варпа организован как вектор на поток,
где для каждого потока хранится отображение типа регистра
на вектор значений:

\begin{lstlisting}[language=C++, caption={Структура регистрового файла}]
using RegisterContext = std::vector<uint64_t>;
using ThreadContext   = std::unordered_map<Ptx::registerType,
                                           RegisterContext>;
// thread_regs[thread_id][registerType][slot]
std::vector<ThreadContext> thread_regs;
\end{lstlisting}

Ключ \texttt{registerType}~--- перечисление \texttt{\{R, Rd, F, Fd, P, ...\}};
вектор значений индексируется номером регистра из PTX-операнда.
Все значения хранятся в \texttt{uint64\_t}: покрывает \texttt{f64}
без потерь; значения \texttt{f32} занимают младшие 32~бита.
Преобразование выполняется через \texttt{std::bit\_cast} без нарушения
правила strict aliasing (C++20).

Размер внутренних векторов определяется директивами \texttt{.reg}
в PTX: исполнитель выделяет ровно столько слотов, сколько объявлено,
без избыточного выделения памяти.

\subsubsection{Вспомогательные шаблоны \texttt{reg\_cast<T>} и \texttt{to\_u64<T>}}

Все исполнители инструкций работают с регистровым файлом через два
шаблонных примитива, определённых в \texttt{executors.tpp}:

\begin{lstlisting}[language=C++, caption={Примитивы доступа к регистрам}]
template<typename T>
T reg_cast(uint64_t raw) {
    if constexpr (std::is_integral_v<T>)
        return static_cast<T>(raw);
    else {
        T val;
        std::memcpy(&val, &raw, sizeof(T));
        return val;
    }
}

template<typename T>
uint64_t to_u64(T val) {
    if constexpr (std::is_integral_v<T>)
        return static_cast<uint64_t>(
                   static_cast<std::make_unsigned_t<T>>(val));
    else {
        uint64_t raw = 0;
        std::memcpy(&raw, &val, sizeof(T));
        return raw;
    }
}
\end{lstlisting}

Целые числа преобразуются через \texttt{static\_cast} через беззнаковый
промежуточный тип~--- бит-паттерн знаковых значений в \texttt{uint64\_t}
сохраняется. Числа с плавающей запятой сериализуются через
\texttt{memcpy}~--- это единственный способ избежать неопределённого поведения
по правилу строгого псевдонимирования (strict aliasing) в C++.

\subsubsection{Вычисление эффективного адреса}

Эффективный адрес потока $t$ для операнда \texttt{[\%rd5+8]}:
\begin{equation}
  ea(t) = R[t][\mathtt{u64}][5] + 8,
\end{equation}
где \texttt{imm\_offset~= 8} извлекается парсером из адресного выражения.
Вычисленный адрес $ea(t)$ передаётся в делегирующий метод
\texttt{WarpContext}, который маршрутизирует обращение в нужное
пространство памяти на основе функции $\mathit{space}(ea(t))$.

\subsection{Исполнители арифметических инструкций}

\subsubsection{Целочисленная арифметика: формальная операционная семантика}

Правила Structural Operational Semantics (SOS) для ключевых целочисленных инструкций:

\textbf{\texttt{add.u32 \%rd, \%ra, \%rb}:}
\begin{equation}
  R'[t][\mathtt{u32}][d] =
  \bigl(R[t][\mathtt{u32}][a] + R[t][\mathtt{u32}][b]\bigr) \bmod 2^{32},
  \quad \forall t \in \activeset(M).
\end{equation}

\textbf{\texttt{mul.hi.u32 \%rd, \%ra, \%rb}:}
\begin{equation}
  R'[t][\mathtt{u32}][d] =
  \left\lfloor \frac{R[t][\mathtt{u32}][a] \cdot R[t][\mathtt{u32}][b]}{2^{32}} \right\rfloor,
  \quad \forall t \in \activeset(M).
\end{equation}

\textbf{\texttt{div.s32 \%rd, \%ra, \%rb}:}
\begin{equation}
  R'[t][\mathtt{s32}][d] =
  \mathrm{trunc}\!\left(\frac{R[t][\mathtt{s32}][a]}{R[t][\mathtt{s32}][b]}\right),
  \quad \forall t \in \activeset(M),\; R[t][\mathtt{s32}][b] \neq 0.
\end{equation}

Граничные случаи: деление на ноль~--- по спецификации PTX возвращает~0;
переполнение \texttt{INT\_MIN / -1}~--- возвращает \texttt{INT\_MIN}
(неопределённое поведение в стандарте C++, обрабатывается явно).

\subsubsection{Арифметика с плавающей запятой: IEEE~754 и режимы округления}

PTX определяет четыре модификатора округления, соответствующие
режимам IEEE~754-2008~\cite{IEEE754-2008}:

\begin{table}[h]
  \caption{Режимы округления PTX и их соответствие IEEE~754}
  \label{tab:rounding}
  \centering
  \begin{tabular}{lll}
    \toprule
    \textbf{Суффикс PTX} & \textbf{Режим IEEE 754} & \textbf{Описание} \\
    \midrule
    \texttt{.rn} & roundTiesToEven  & к ближайшему чётному \\
    \texttt{.rz} & roundTowardZero  & усечение \\
    \texttt{.ru} & roundTowardPositive & к $+\infty$ \\
    \texttt{.rd} & roundTowardNegative & к $-\infty$ \\
    \bottomrule
  \end{tabular}
\end{table}

Модификаторы округления (\texttt{.rn}, \texttt{.rz}, \texttt{.ru},
\texttt{.rd}) разбираются парсером и сохраняются в объекте инструкции.
В текущей реализации эмулятор не переключает режим округления хост-CPU
через \texttt{fesetround}: операции с плавающей запятой выполняются
с режимом \texttt{FE\_TONEAREST} по умолчанию. Это известное
ограничение, которое не влияет на корректность результата для
инструкций, не чувствительных к последнему биту (\textit{faithful
rounding}).

\subsubsection{Инструкции преобразования типов}

\textbf{\texttt{cvt.rni.s32.f32 \%rd, \%fa}:}
преобразование \texttt{f32} $\to$ \texttt{s32} с округлением к ближайшему целому:
\begin{equation}
  R'[t][\mathtt{s32}][d] = \left\lfloor R[t][\mathtt{f32}][a] + 0.5 \right\rfloor,
  \quad \forall t \in \activeset(M).
\end{equation}

\textbf{\texttt{cvt.rz.f32.s64 \%fd, \%ra}:}
преобразование \texttt{s64} $\to$ \texttt{f32} с усечением дробной части:
\begin{equation}
  R'[t][\mathtt{f32}][d] = \mathrm{round}_z\bigl(R[t][\mathtt{s64}][a]\bigr),
  \quad \forall t \in \activeset(M).
\end{equation}

\subsection{Исполнители инструкций работы с памятью}

\subsubsection{\texttt{ld.global} / \texttt{ld.shared}: чтение}

Исполнитель \texttt{ld} реализован как шаблонный метод \\
\texttt{ldInstruction::ExecuteThread<dataType~Data>}, параметризованный
типом данных. Тип \texttt{ptx\_native\_t<Data>}~--- псевдоним,
отображающий перечислитель \texttt{dataType} на соответствующий
C++-тип (\texttt{F32}~$\to$~\texttt{float}, \texttt{S64}~$\to$~\texttt{int64\_t} и т.д.):

\begin{lstlisting}[language=C++, caption={ldInstruction::ExecuteThread (executors.tpp)}]
template<dataType Data>
void ldInstruction::ExecuteThread(
        uint32_t lid, std::shared_ptr<WarpContext>& wc)
{
    using T = ptx_native_t<Data>;
    uintptr_t addr = 0;
    if (addr_.reg.type != registerType::UNDEFINED) {
        addr = static_cast<uintptr_t>(
            wc->thread_regs[lid][addr_.reg.type][addr_.reg.reg_id]);
        addr += static_cast<ptrdiff_t>(addr_.imm);
    }
    if (space_ == ldspaceQl::Shared) {
        auto* base = static_cast<uint8_t*>(wc->getSharedBase());
        addr = reinterpret_cast<uintptr_t>(base + addr);
    }
    T val = T{};
    std::memcpy(&val, reinterpret_cast<const void*>(addr), sizeof(T));
    wc->thread_regs[lid][dst_.type][dst_.reg_id] = to_u64<T>(val);
}
\end{lstlisting}

Метод вызывается по одному разу для каждого активного потока варпа.
Переключение между пространствами памяти (\texttt{global}/\texttt{shared})
производится внутри метода по полю \texttt{space\_}, разобранному
парсером. Данные копируются через \texttt{memcpy}~---
без нарушения strict aliasing и без промежуточного приведения указателей.

\subsubsection{\texttt{st.global} / \texttt{st.shared}: запись с учётом маски}

Для каждого потока $t \in \activeset(M)$:
$*\mathit{ptr}(ea(t)) \leftarrow R[t][\tau][\mathit{src}]$.
Потоки с $M[i]=0$ пропускают запись~--- это семантически обязательно:
неактивные потоки не должны изменять состояние памяти даже при
вычислении адреса.

\subsubsection{\texttt{setp} / \texttt{selp}: операции с предикатами}

Для \texttt{setp.lt.f32 \%p, \%fa, \%fb}:
\begin{equation}
  P'[t] = \bigl(R[t][\mathtt{f32}][a] < R[t][\mathtt{f32}][b]\bigr),
  \quad \forall t \in \activeset(M).
\end{equation}

Результат сравнения (\texttt{true}/\texttt{false}) записывается в один
предикатный регистр \texttt{dst\_} как значение \texttt{1}/\texttt{0}
типа \texttt{uint64\_t}.

\subsection{Исполнители инструкций управления потоком}

\subsubsection{Безусловный переход \texttt{bra.uni \$L}}

\begin{equation}
  \eval{\mathtt{bra.uni}\;\$L}{\sigma}
  \reduce \sigma[\mathrm{PC} \leftarrow \mathit{offset}(\$L)].
\end{equation}

Суффикс \texttt{.uni} (uniform) гарантирует конвергентность.
В реализации: \texttt{warp.pc = warp.GetBasicBlockPc("\$L")}.
Стек дивергенции не модифицируется.

\subsubsection{Условный переход \texttt{@\%p bra \$L}}

Полная семантика с разбором по случаям (подробно в Главе~6):

$M_{\mathrm{taken}} = M \wedge P$, $M_{\mathrm{not}} = M \wedge \neg P$.

\begin{itemize}
  \item \textbf{Случай 1} ($M_{\mathrm{not}} = 0$):
        $\mathrm{PC}' = \mathit{offset}(\$L)$.
  \item \textbf{Случай 2} ($M_{\mathrm{taken}} = 0$):
        $\mathrm{PC}' = \mathrm{PC} + 1$.
  \item \textbf{Случай 3}: Вызов PDOM\_Diverge (см.~Главу~6).
\end{itemize}

Инструкция \texttt{ret} завершает исполнение варпа: она сбрасывает
устаревшие фреймы PDOM-стека дивергенции и устанавливает
\texttt{pc = EOC} (End-of-Code), сигнализируя циклу исполнения
об окончании:

\begin{lstlisting}[language=C++, caption={retInstruction::ExecuteBranch (executors.tpp)}]
// Discard stale convergence frames from guard-clause paths.
while (!wc->execution_stack.empty()
       && wc->execution_stack.top().is_convergence)
    wc->execution_stack.pop();

if (wc->execution_stack.empty())
    wc->pc = WarpContext::EOC;
else {
    const auto frame = wc->execution_stack.top();
    wc->execution_stack.pop();
    wc->pc             = frame.pc;
    wc->execution_mask = frame.mask;
}
\end{lstlisting}

\subsubsection{\texttt{bar.sync}: делегирование барьерной синхронизации}

Исполнитель \texttt{bar.sync~N} вызывает
\texttt{warp.syncBarrier()}, реализованный в компонентах управления
блоком (Часть~2).

Математически: все записи \texttt{st.shared} до \texttt{bar.sync}
гарантированно видимы всем потокам после него~---
барьер обеспечивает полную видимость памяти
(happens-before по отношению к записям разделяемой памяти).
\newpage
